\begin{papercard}{krippendorff2019}
\cf{Stages} \stagetag{C} \stagetag{I} (also \stagetag{S} \stagetag{R}). The standard textbook of content analysis (4th ed.). This card covers the design logic and the data-making chapters (4 Logic of designs, 5 Unitizing, 6 Sampling, 7 Recording/Coding); the reliability chapter (12) and the practical guide (14) are below. \emph{Note:} page numbers refer to PDF pages of the appended copy, i.e.\ book page + 15 (e.g.\ PDF page 101 = book page 86). The appended full text is a private copy, kept out of the repository.

\cf{Chapter 4: logic of designs} Data are ``made, not found'', so researchers must say how they generated them \pc{krippendorff2019}{101}{Hence data are made, not found, and researchers are obligated to say how they generated their data}. A design's logic has two qualities: \emph{efficiency} and \emph{evenhandedness}. For replicability the account must be complete enough to serve as instructions for coders and critics \pc{krippendorff2019}{101}{the researcher’s descriptive account of the analysis must be complete enough to serve as a set of instructions to coders}. Six components, each with a descriptive and an operational state: unitizing, sampling, recording/coding, data reduction, abductive inference, narration. The first four form \emph{data making} \pc{krippendorff2019}{103}{Together, the first four components constitute what may be summarily called data making}. The process is recursive in \emph{development}, but the \emph{execution} of the finished design should be fixed and replicable. Analysts should ``resist the violence'' of poor instructions during development \pc{krippendorff2019}{106}{a sensible analyst should “resist the violence” that poor instructions can inflict on the texts}. Quantification is not a defining criterion, and qualitative and quantitative analysis differ in justification (explicitness vs.\ contextual appropriateness), not in kind \pc{krippendorff2019}{106}{Using numbers instead of verbal categories or counting instead of listing quotes is merely convenient}. Three preparatory designs yield an analytical construct: operationalising knowledge of the context, testing constructs as hypotheses, and developing constructs by trial and error with feedback \pc{krippendorff2019}{109}{Below, I discuss three ways of establishing analytical constructs}.

\cf{Chapter 5: unitizing} Units are not givens, since they emerge in reading and depend on the analyst's ability to see meaningful conceptual breaks \pc{krippendorff2019}{118}{units should not be considered givens}. Three kinds: \emph{sampling units} (selected for inclusion; must not leave out relevant information or bias across connections), \emph{recording/coding units} (separately described; contained in or equal to sampling units; choose the smallest unit that carries the needed information) and \emph{context units} (limit the information consulted; not counted; may overlap) \pc{krippendorff2019}{120}{Context units are units of textual matter that set limits on the information to be considered}. ``Context units should be as large as is meaningful (validity) and as small as is feasible (reliability)'' \pc{krippendorff2019}{121}{The best content analyses define their context units as large as is meaningful}. Five ways to define units: physical, syntactical, categorial, propositional, thematic \pc{krippendorff2019}{122}{Content analysts identify units according to one or more of five kinds of distinctions}. Trade-off: thematic units are rich but hard to code reliably, simpler ones are reliable but less productive \pc{krippendorff2019}{129}{Thematic distinctions are rich in information and potentially very productive}. The reliability of unitizing itself must be assessed (section 12.6) \pc{krippendorff2019}{129}{Calculating the reliability of unitizing may not be simple}.

\cf{Chapter 6: sampling} Sampling theory for individuals does not map onto texts: units are not independent, sampling and counted units differ, texts are unequally informative, and analysts must sample one population (texts) with a view to another (answers) \pc{krippendorff2019}{131}{Sampling theory, as outlined above, does not, however, fully map out the sampling problems that content analysts face}. Techniques: random, systematic (beware of cyclic bias), stratified, varying-probability, cluster, snowball, \emph{relevance (purposive)}, census, convenience \pc{krippendorff2019}{137}{Relevance sampling, in contrast, aims at selecting all textual units that contribute to answering given research questions}. Internet sampling needs narrow time windows because texts change while being coded \pc{krippendorff2019}{132}{The changing nature of texts on the Internet calls for sampling plans with narrow time windows}. Convenience samples leave undecided how and why data entered the sample and may deceive the analyst \pc{krippendorff2019}{139}{Convenience samples present content analysts with the potential problem of having to undo or compensate for the biases}. Sample size: if relevant units are rare, the sample must be larger (Table 6.1: p = .001 at 0.05 needs 2,995), experiments against a benchmark, or the \emph{split-half} technique \pc{krippendorff2019}{142}{The split-half technique calls for analysts to divide a sample randomly into two parts of equal size}.

\cf{Chapter 7: recording/coding} Recording translates transient or unstructured phenomena into a data language, and coding is recording by observer-independent rules. Recording instructions must cover coder qualifications, training, the data language (syntax and semantics) and the records \pc{krippendorff2019}{144}{The qualifications that coders need to have}. The codebook must survive with the data, because without it records are ``computable but no longer interpretable'' \pc{krippendorff2019}{145}{the records that a study has generated are reduced to nothing more than a collection of meaningless marks or numbers}. \emph{Coder qualifications}: cognitive ability, shared background (describe it so others can recruit similar coders) and \emph{frequency} of that qualification in the population, so the study is replicable \pc{krippendorff2019}{146}{Any researcher who claims to be the only one who is capable of reading a text correctly in fact denies the possibility of replicating the research}. \emph{Training}: reliability tests during development until the requirement is met. Warning: agreement between coders who co-developed the scheme is inflated by tacit consensus, so test the final instructions with a \emph{fresh set of coders} working independently from the instructions alone \pc{krippendorff2019}{148}{the final form of the coding instructions should be tested for their replicability with a fresh set of coders}. Self-applied categories are ``notoriously unreliable'' \pc{krippendorff2019}{149}{self-applied recording instructions are notoriously unreliable}. Crowdcoding (mTurk) measures agreement on transparent instructions, but screening, test questions and shared-ability limits apply \pc{krippendorff2019}{151}{agreement observed among crowdworkers is not so polluted}.

\cf{Data-language semantics (7.5)} Categories must be \emph{mutually exclusive and exhaustive}. Catch-all categories (``other'') add little, and ``undecidable'' categories reveal unclear concepts rather than fix them \pc{krippendorff2019}{154}{The well-intended practice of adding such categories as “undecidable,” “ambiguous,” or “applicable to two or more categories”}. Devices: verbal designations (plain, widely understood names; extend binary distinctions), extensional lists, examples (prototype-near examples of what to include and exclude), \textbf{decision schemes} (sequenced decisions; people hold about 7$\pm$2 alternatives; clarifies logical levels; Carletta's conversational-move tree) \pc{krippendorff2019}{157}{As a rule of thumb, humans cannot keep the meanings of more than seven (plus or minus two) alternatives}, magnitudes and scales (not all text is scalable, and a midpoint confuses ``uncodable'' with ``balanced''; unipolar 0--6 scales as an alternative) \pc{krippendorff2019}{159}{if coders are instructed to mark the midpoint of a scale whenever they encounter something that is not codable}, simulation of hypothesis testing (affirmed, implicitly affirmed, irrelevant, implicitly denied, denied), simulation of interviewing, and constructs for closure \pc{krippendorff2019}{161}{Such a testing of mutually exclusive propositions (hypotheses)}.

\cf{Records (7.6)} One record per recording unit with the same variables. Records need \emph{administrative information}: project, version of instructions, kind of text (pilot or main), state (completed, verified), coder identity, and the unit's serial number \pc{krippendorff2019}{166}{Administrative information guides the handling of data}. Multilevel units need a master file, cross-references or fully redundant records. Design errors to avoid: overusing numbers, inconsistent codes for ``other/not applicable'' across variables, free-text comments instead of codes, and recording formats that invite illegal entries (prefer ``check all that apply'') \pc{krippendorff2019}{169}{A third source of errors is the adding of comments to records}, \pc{krippendorff2019}{170}{The most common source of coding errors is evident when the recording instructions are not well spelled out}. Recommendation: present coders with the variable, its mutually exclusive options and the six MacQueen codebook fields at the point of entry, with elaborate instructions hidden but immediately available \pc{krippendorff2019}{170}{The marker (code) that actually is recorded}.

\cf{Chapter 12: why reliability and what kinds} Reliable data are ``data that remain constant throughout variations in the measuring process'' \pc{krippendorff2019}{292}{Reliable data, by definition, are data that remain constant throughout variations in the measuring process}. Two operationalisations: measurement-theory (independent of the measuring event) and interpretivist (a community concurs on readings). Reliability is \emph{necessary but not sufficient} for validity, and chasing reliability can cost validity \pc{krippendorff2019}{294}{Thus, reliability is a necessary, but not a sufficient, condition for validity}, \pc{krippendorff2019}{294}{In the pursuit of high reliability, validity tends to get lost}. Table 12.1 gives three \emph{types}: \textbf{stability} (test--retest; ``the weakest reliability measure in content analysis'' because it rewards stable idiosyncrasy), \textbf{replicability} (test--test, independent coders), and \textbf{accuracy or surrogacy} (against a standard; the strongest, but standards are rarely available) \pc{krippendorff2019}{296}{Stability is the weakest reliability measure in content analysis}, \pc{krippendorff2019}{296}{Replicability is a measure of the degree to which a process can be reproduced by different analysts}. The \emph{split-half technique} is not a reliability test \pc{krippendorff2019}{297}{the split-half technique is not a suitable reliability test}.

\cf{Conditions for generating reliability data} Communicable instructions (including accounts of coder training), communicable criteria for selecting coders from a population of equally capable people, and \emph{independent} observers \pc{krippendorff2019}{298}{It must ensure that the observers who generate the reliability data work independently of each other}. Practices that fake reliability: coders reaching decisions by consensus, which ``afford no reliability test'' \pc{krippendorff2019}{299}{the data that are generated by consensus afford no reliability test}; coders consulting each other and silently extending instructions; relying on a single expert or on like-minded friends; and discarding units on which coders disagree \pc{krippendorff2019}{300}{Sometimes content analysts accept as data only those units of analysis on which observers achieve perfect agreement}. Recommended: three or more independent observers, then reconcile by rule or deliberation. ``The only publishable reliability is the one measured \emph{before} the reconciliation of disagreements'' \pc{krippendorff2019}{300}{The only publishable reliability is the one measured before the reconciliation of disagreements}.

\cf{Krippendorff's $\alpha$} Reliability data are an observers-by-units matrix (canonical form). Properties: corrects for chance, independent of the number of values, any number of observers, unequal sample sizes, all metrics (nominal, ordinal, interval, ratio), missing data \pc{krippendorff2019}{306}{It is applicable despite missing data caused by unequal numbers of observers attending to unequal sets of recording units}. General form $\alpha=1-D_o/D_e$ (observed vs.\ expected disagreement): 1 means perfect reliability and 0 means chance \pc{krippendorff2019}{306}{α assesses the deviation from perfect reliability by the proportion of observed to expected disagreement}. Negative values arise from small samples or \emph{systematic} disagreements \pc{krippendorff2019}{307}{Systematic disagreements are more serious}. Not Cronbach's $\alpha$ \pc{krippendorff2019}{306}{α-agreement should not be equated with Cronbach’s (1951) alpha}. Software: KALPHA macros for SPSS/SAS (Hayes \& Krippendorff 2007); check results by hand because flawed packages exist \pc{krippendorff2019}{308}{Useful SAS and SPSS macros for computing}. Variants: $c\alpha$ for coding predefined units, $\alpha_u$/$_{|u}\alpha$/$_{cu}\alpha$/$_{(k)u}\alpha$ for \emph{unitizing} continua (12.6), $\alpha$ for multi-valued coding (12.7), and crowd/accuracy variants (12.4).

\cf{Contested coefficients (12.3.5)} Only 56\,\% of content analyses reported reliability in one survey, and 35\,\% used percent agreement. \emph{Correlation coefficients} measure covariation, not agreement, and are ``seriously misleading'' \pc{krippendorff2019}{331}{Although correlation coefficients have been used as reliability measures, even recommended}. \emph{Percent agreement} ignores chance; Bennett's $S$ depends on the number of categories and can be inflated by unused categories. Scott's $\pi$ and Fleiss' $K$ estimate expected agreement from \emph{pooled} category proportions. Cohen's $\kappa$ uses each coder's marginals, so it ``overestimates reliability whenever coders have unequal preferences for coding categories'' \pc{krippendorff2019}{333}{κ overestimates reliability whenever coders have unequal preferences for coding categories}, and it is a ``hybrid'' whose expected agreement resembles a $\chi^2$ association statistic \pc{krippendorff2019}{335}{Cohen’s κ, by contrast, reveals itself as a hybrid coefficient}. Gwet's AC1 is hard to interpret and yields higher values without telling whether imperfect data are usable \pc{krippendorff2019}{335}{Gwet’s AC1 is even more difficult to interpret}. Cronbach's $\alpha$ ``must not be construed as an agreement measure'' \pc{krippendorff2019}{336}{It belongs to the family of correlation coefficients and must not be construed as an agreement measure}. Only $\alpha$ handles more than two coders, so the common two-coder practice is ``woefully inadequate'' \pc{krippendorff2019}{337}{has given rise to the dominant practice of testing coding reliability with just two coders}. Calling coefficients ``liberal'' or ``conservative'' misses that they respond to different things \pc{krippendorff2019}{337}{the choice among coefficients is attitudinal and has to do with the numbers researchers feel comfortable reporting}. Numerical demonstration: two tables with the same 46\,\% agreement give $\kappa$ 0.186 and 0.258 but $\pi=\alpha=0.186$.

\cf{Insufficient variation, significance, sample size} If almost all units get the same value, $\alpha$ is near 0 even at 88\,\% agreement, because there is no information about reliability. Rare categories need larger samples \pc{krippendorff2019}{367}{Time and time again, statistical novices find this condition difficult to accept}, \pc{krippendorff2019}{367}{Data without variation not only fail to pass reliability tests}. Krippendorff discusses Feinstein \& Cicchetti's paradox and Gwet's remedy but prefers to assess variation separately \pc{krippendorff2019}{367}{Feinstein and Cicchetti (1990) and Cicchetti and Feinstein (1990) have discussed the problem of low variation}. \emph{Significance}: rejecting the null of chance agreement is the wrong test, since the aim is to show data are close to perfect \pc{krippendorff2019}{368}{A common mistake is to accept or reject data as reliable when the null hypothesis that agreement occurs by chance can be rejected}. Use a \emph{bootstrapped} confidence interval and the probability $q$ of failing to reach $\alpha_{\min}$ \pc{krippendorff2019}{368}{Bootstrapping involves drawing several thousand subsamples from the available reliability data}. \emph{Sample size}: each category should yield at least five chance agreements as a rule of thumb, and Table 12.3 gives required numbers of values (e.g.\ 241 values for $P_{\min}=.125$, $\alpha_{\min}=.800$, $p=.05$, i.e.\ 121 units for two coders) \pc{krippendorff2019}{369}{As a rule of thumb, each category of units should occur often enough to yield at least five agreements by chance}, \pc{krippendorff2019}{371}{the table suggests a minimum number of 241 values}.

\cf{Standards (12.8.4)} Aim at $\alpha=1$; rely only on variables with $\alpha\ge.800$; use $.667\le\alpha<.800$ only for tentative conclusions \pc{krippendorff2019}{371}{Rely only on variables with reliabilities α ≥ 0.800}, \pc{krippendorff2019}{371}{Consider variables with reliabilities between α = 0.667 and α = 0.800 only for drawing tentative conclusions}. A reference point: coders using Dutch words they did not understand still reached $c\alpha=0.44$ \pc{krippendorff2019}{371}{the words must have invoked some consistent associations with perceived personality characteristics}. In confidence-interval terms: do not accept data whose interval reaches below $\alpha_{\min}$, and keep $q$ small (e.g.\ .05). Standards ``should always be related to the validity requirements'' and the cost of wrong conclusions \pc{krippendorff2019}{372}{I recommend such levels with considerable hesitation}. Verbal labels such as Landis \& Koch's are ``largely misleading'' \pc{krippendorff2019}{372}{I am suggesting that benchmarks like these are largely misleading}. Never average the reliabilities of different variables: ``the lowest $\alpha$ among them is the reliability of all'' \pc{krippendorff2019}{373}{It is a serious mistake to average the reliabilities of the variables of a complex coding instrument}, \pc{krippendorff2019}{374}{For data with many variables, the lowest α among them is the reliability of all}. Ultimate test: replicate the data-making for all or a large share of the data, generate versions by permuting disagreements, and check that conclusions agree \pc{krippendorff2019}{374}{Replicate the data-making process for all or for a large enough proportion of data being analyzed}.

\cf{Chapter 14: practical guide} Three entries: \emph{text-driven} (``fishing expeditions''), \emph{problem-driven} and \emph{method-driven} \pc{krippendorff2019}{399}{Problem-driven content analyses (14.1.2) are motivated by epistemic questions}. A research question concerns unobserved phenomena, has several possible answers and at least two ways of selecting among them \pc{krippendorff2019}{402}{They concern currently unobserved phenomena in the problematized context of available texts}, and counting alone is not a research question \pc{krippendorff2019}{402}{Counting is an operation performed on a body of text}. Specify before collecting data what could validate or invalidate the result \pc{krippendorff2019}{403}{content analysts should be specific about what could validate or invalidate their results before collecting data}. Nine problem-driven steps: question, stable correlations, locate texts, sampling units, sample, coding categories and instructions, analytical procedures, reliability standards, resources. For instructions: a list of coder qualifications, training materials, operational unit definitions, the data language with examples of what \emph{not} to code, and the record forms. Pretest and revise until reliability is adequate \pc{krippendorff2019}{409}{Typically, before these instructions are applied by several coders and to a large body of text, the analysts need to pretest them}. Prefer categories proven elsewhere, since bespoke categories make results non-comparable \pc{krippendorff2019}{409}{Unfortunately, many content analysts use categories that are uniquely tailored to available texts}, \pc{krippendorff2019}{410}{If existing coding categories have proven reliable and if they lead to answers to the research question under consideration, there is no need to invent a new scheme}, and draw categories from theory or official taxonomies \pc{krippendorff2019}{410}{A third strategy is to draw from available literature on or theories of the context of the analysis}. Ethics: websites and public documents are ethically neutral, but writing not meant to be public can be harmed even under pseudonyms \pc{krippendorff2019}{407}{However, content analysts can easily and unintentionally violate the privacy of individuals}. The most frequent disruption in application is failure to meet reliability standards \pc{krippendorff2019}{418}{the most frequent disruptions stem from the inability to meet accepted standards of reliability}. A proposal and a report outline (14.2, 14.4.1) include: units, sampling, data language, reliability and validity, techniques, appendices with the recording instructions and per-variable reliabilities \pc{krippendorff2019}{421}{The recording/coding process: whether built into computer dictionaries or search queries}, \pc{krippendorff2019}{421}{Appendices containing materials that enable interested readers to read beyond the report, such as the recording instructions}.

\cf{Use for this thesis} The reference text for the \emph{coding-reliability design}. (1) Unit definitions: state the sampling unit (app, store, version), recording unit (walkthrough step or observed element) and context unit (screen, flow), and report the reliability of unitizing if units are judged. (2) The codebook is a recording instruction: include coder qualifications, training, data language with mutually exclusive and exhaustive categories (avoid ``undecidable''), decision schemes for confusable codes, and examples of what \emph{not} to code. (3) Test the \emph{final} instructions with a coder who did not co-develop them, working independently, and report the \emph{pre-reconciliation} reliability. (4) Report $\alpha$ (with bootstrapped confidence interval and the sample-size logic), with the thresholds $.800$ / $.667$ treated as conventions to justify, plus per-variable values and no averaging across variables. (5) Anticipate prevalence skew: low-variation codes need targeted sampling, and the trade-off against $\kappa$-type statistics (see Feinstein \& Cicchetti, Gwet, Hallgren) must be decided and logged. (6) Record administrative information (instruction version, coder, state) in every record, which the registry and capture log can provide. Caution: this is a 2019 \emph{textbook} copy; its criticisms of $\kappa$ and AC1 are the author's position and are contested, so cite the contested points as positions and not as consensus.
\end{papercard}
